\begin{chapintro}
検査のレビューの問いは、対象、対象検査、参照基準、標的疾患の 4 つと、検査を使う場面と役割で決まります。検索では研究デザインの絞り込み（フィルター）をかけず、広く拾ってから人の目で選別します。架空の PCT の抄録 7 件で、組み入れるかどうかの判断を練習します。不完全な参照基準を扱う潜在クラスモデルの紹介は、発展の囲みに置いてあり、初読では飛ばせます。
\end{chapintro}

\section{問いを立てる}
\label{sec:11-1}

介入のメタ解析の問いを PICO（対象、介入、比較、アウトカム）で立てたように、検査のメタ解析の問いは次の要素で組み立てます\kksee{sec:02-1}。
\begin{description}
\item[対象] どんな人に検査を行うのかを決めます。例題では、発熱して救急外来を受診した成人、または ICU で発熱した成人です。
\item[対象検査] 精度を知りたい検査とカットオフです。例題では血清または血漿の PCT で、カットオフは研究ごとに異なってよいとします。
\item[参照基準] 正解を決める方法です。例題では血液培養で、汚染菌と判断された陽性は陰性として扱います。
\item[標的疾患] 見分けたい状態です。例題では菌血症です。
\end{description}

この 4 つに加えて、検査を臨床のどの場面で、どんな役割で使うのかを決めておきます。役割には、既存の検査の置き換え（今の検査の代わりに使う）、トリアージ（後の検査を受ける人を絞り込むために先に使う）、追加（今の検査に上乗せする）があります。例えば PCT を「血液培養を採るかどうかを決めるトリアージ」として使うなら、見逃し（偽陰性）を少なくすることが求められ、感度を重く見ます。「抗菌薬を始めるかどうかの決め手」として使うなら、偽陽性による不要な投与も問題になります。役割によって、結果のどの数字を重く見るかと、どの研究がその問いに合っているかが変わります。

例題の参照基準の血液培養は、完全な正解ではありません。採血の量、採るセットの数、抗菌薬を投与した後かどうかで、菌血症のある人が陽性になる割合が変わります。菌血症があっても培養が陰性になった人は、2×2 表では菌血症なしに数えられ、PCT が高ければ偽陽性になるので、PCT の特異度は低く見えることがあります。

\begin{hattenbox}
不完全な参照基準を扱う方法として潜在クラスモデルがあり、Cochrane DTA ハンドブックの 9.5.1 に書かれていますが、本書では扱いません。
\end{hattenbox}

レビューでは、参照基準の限界を抽出フォームに記録し（\kkref{sec:12-1}）、考察で述べます。

例題の問いを 1 文で書くと、「発熱した成人（救急外来または ICU）で、PCT は血液培養で確認された菌血症をどれだけ正しく見分けるか」になります。この問いは、\kkref{sec:02-4}で説明したようにプロトコルに書き、PROSPERO に登録してから検索を始めます。プロトコルには、カットオフが複数報告されている研究からどの値を取るか、判定不能の結果をどう扱うかといった、データ抽出の規則も先に書いておきます（\kkref{sec:12-3}）。

\section{検索の組み立て}
\label{sec:11-2}

検索の考え方は\kkref{sec:03-2}と同じで、概念ごとに同義語を OR でつなぎ、概念どうしを AND でつなぎます。検査のレビューで使う概念は、多くの場合「対象検査」と「標的疾患」の 2 つだけです。PubMed での検索式の例を示します。

\begin{rcode}
("Procalcitonin"[Mesh] OR procalcitonin[tiab] OR PCT[tiab])
AND
("Bacteremia"[Mesh] OR bacteremia[tiab] OR bacteraemia[tiab]
 OR "bloodstream infection"[tiab] OR "blood culture"[tiab])
\end{rcode}

この式は R のコードではなく、PubMed の検索欄に入れる文字列です。介入の RCT を探すときは「ランダム化比較試験」に絞り込む検索フィルターを付けることがありますが、診断精度研究ではこのような研究デザインのフィルターを付けません。診断精度研究は、論文の題や抄録に「感度」「特異度」と書かれていないことが多く、索引語（MeSH）の付け方もそろっていないので、フィルターを付けると必要な研究を取りこぼします。Cochrane DTA ハンドブックも、フィルターだけで絞り込むことを勧めていません。その代わり、選別する抄録の数は増えます。

データベースは PubMed（MEDLINE）、Embase、Cochrane Library を基本にし、日本の研究を拾うなら医学中央雑誌も加えます。検査の研究は学会抄録や企業の技術資料として発表されることもあるので、灰色文献をどこまで探すかもプロトコルに書いておきます。検索式とその実行日は、PRISMA-DTA に従って報告するため、データベースごとに記録します。

\section{抄録の選別の手順}
\label{sec:11-3}

検索で集めた記録は、まず文献管理ソフトで重複を取り除きます。そのうえで、2 人のレビュー著者が独立に、題と抄録だけを読んで、組み入れ基準を満たす可能性があるかを判定します。判定は 3 通りで、「明らかに満たさない」なら除外し、「満たす」または「分からない」なら本文を取り寄せて確認します。抄録の段階では、迷ったら残すのが原則です。抄録に書かれていないことは本文に書かれていることが多く、抄録の段階で除外した研究は二度と見直されないからです。詳しくは\kkref{sec:03-3}を読み直してください。

2 人の判定が食い違ったら、話し合って決め、決まらなければ 3 人目が判定します。本番の前に 50 件ほどを 2 人で試しに判定し、判定の基準をそろえておくと、食い違いが減ります。2 人の判定の一致の度合いを κ 係数で報告することもあります。

\section{抄録の選別の演習}
\label{sec:11-4}

架空の検索で見つかった抄録のうち 7 件を、短く要約して示します。どれも架空のもので、実在の論文ではありません。組み入れ基準は、先の問いの 4 要素を満たし、2×2 表が作れることとします。読者も判定してから、次の節の答えと見比べてください。第 1 部の\kkref{sec:03-4}と同じ形式の演習です。

\begin{description}
\item[抄録 A（Lopez 2011）] スペインの救急外来で、発熱した成人 482 人を連続して登録した。受診時に PCT を測定し、全員に血液培養を行った。血液培養陽性は 62 人だった。カットオフ 0.5 ng/mL での感度と特異度を報告した。
\item[抄録 B（Mori 2013）] 日本の ICU で、新たに発熱した成人 178 人を前向きに登録し、発熱時の PCT と血液培養を比べた。血液培養陽性は 38 人で、カットオフ 0.5 ng/mL での診断精度を示した。
\item[抄録 C] 血液培養で菌血症と確定した入院患者 60 人と、健康診断を受けた健康な成人 60 人の PCT を比べた。カットオフ 0.5 ng/mL で感度 0.92、特異度 0.98 だった。
\item[抄録 D] 救急外来の発熱患者 640 人で、PCT と CRP の菌血症に対する AUC を比べた。PCT の AUC は 0.79、CRP は 0.66 だった。感度と特異度は抄録に書かれていない。
\item[抄録 E] 小児救急外来を受診した 3 か月から 15 歳の発熱児 410 人で、PCT による重症細菌感染症の診断精度を調べた。
\item[抄録 F] ICU の成人 230 人で、主治医が臨床経過から判断した「敗血症」の有無を正解として、PCT の診断精度を調べた。
\item[抄録 G] 下気道感染症の外来患者を、PCT の値に基づいて抗菌薬を処方する群と通常診療の群にランダムに割り付け、抗菌薬の処方日数を比べた。
\end{description}

\section{演習の答えと判定の理由}
\label{sec:11-5}

判定と理由を表 11.1 にまとめます。

\begin{kkcenter}
\kkfit{\begin{tabular}{lll}
\toprule
抄録 & 判定 & 理由 \\
\midrule
A & 本文へ & 対象、検査、参照基準、標的疾患がすべて合う \\
B & 本文へ & ICU の発熱患者で、問いの範囲に入る \\
C & 除外 & 症例対照デザインで、対照が健康な人 \\
D & 本文へ & 抄録には 2×2 がないが、本文で作れる可能性がある \\
E & 除外 & 対象が小児で、問いの対象（成人）と違う \\
F & 除外 & 参照基準が主治医の臨床判断で、標的疾患も敗血症 \\
G & 除外 & 精度ではなく、検査を使った治療方針の効果を調べた RCT \\
\bottomrule
\end{tabular}}
\kkcap{表 11.1　架空の抄録 7 件の判定}
\end{kkcenter}

抄録 A と B は迷わず残します。抄録 D は、抄録に感度と特異度が書かれていないので除外したくなりますが、本文の表や補足資料に複数のカットオフでの成績が載っていることがあります。2×2 表が作れるかどうかは本文を見ないと分からないので、本文へ進めます。本文を見ても作れなければ、その時点で「2×2 表が作れない」という理由で除外し、理由を記録します。

抄録 C は、感度 0.92、特異度 0.98 と、ほかの研究よりずっと良い成績を示しています。この成績の良さは、研究の設計から生まれたものです。対照群が健康診断を受けた健康な人なので、PCT はほとんど上がっておらず、特異度は高く出ます。症例群も、入院してから菌血症と確定した、病状のはっきりした患者なので、PCT は高く、感度も高く出ます。実際に PCT を使う場面で見分けたいのは、同じように発熱して受診した人の中の、菌血症の人とそうでない人です。発熱していて PCT がそこそこ上がっているけれども菌血症ではない人（ウイルス感染症や非感染性の炎症）が対照にいないので、見分けにくい部分が研究から抜け落ちています。

このように、研究に含まれる患者の病状の幅（スペクトラム）が実際の使い方と違うことで精度が偏ることを、スペクトラムバイアスと呼びます。症例対照デザインの診断精度研究は精度を高めに見積もりやすく（表 12.1、\kkref{sec:12-5}）、QUADAS-2（\kkref{sec:12-4}）でも、患者選択の領域の信号質問「症例対照デザインを避けたか」に「いいえ」と答えることになり、バイアスのリスクは高いと判定されやすくなります。

症例対照デザインの研究を組み入れ基準で除外するか、組み入れたうえで患者選択を高と判定するかは、プロトコルで事前に決めます。例題では、連続登録またはランダム登録のコホートを組み入れ基準にして、抄録 C のような研究は除外としました。

抄録 E は、対象が小児なので除外します。小児と成人では、PCT の基準範囲も、菌血症の起因菌や頻度も違うので、同じ問いには入れません。小児を含めるかどうかはプロトコルで決めておくことで、データを見てから決めることではありません。

抄録 F は、参照基準が主治医の判断で、標的疾患も菌血症ではなく敗血症なので、問いに合いません。さらに、主治医が PCT の値を知ったうえで敗血症かどうかを判断していれば、正解の決め方が検査の結果に左右されます。参照基準を判定する人が対象検査の結果を知っていることによる偏りを、診断レビューバイアスと呼びます（表 12.1）。QUADAS-2 では参照基準の領域で問います。これとは別に、対象検査の結果そのものが参照基準の一部になっている場合を、取り込みバイアス（incorporation bias）と呼びます。たとえば、診断基準に血清クレアチニンの上昇が含まれる急性腎障害を、血清クレアチニンで見分けられるかを調べる研究では、正解の決め方に検査の値が入っています。PCT の値を診断の基準に含めて敗血症と判定した研究も同じです。抄録 F の記載だけでは、主治医が PCT の値を知っていたのか、診断の根拠に PCT が入っていたのかは分からないので、本文で確かめます。どちらの場合も、PCT と正解が実際より一致して見え、精度は高めに出る恐れがあります。

抄録 G は、PCT の精度ではなく、PCT を使った診療方針の効果を調べた RCT です。これは介入のメタ解析の対象で、検査の精度のレビューには入れません。

\section{フロー図}
\label{sec:11-6}

選別の結果は、PRISMA-DTA に沿ったフロー図で報告します。PRISMA 2020 の書式に合わせ、研究の数と報告の数を分けて数えます（フロー図の見方は\kkref{sec:03-5}で、重複した報告を試験単位にまとめる作業は\kkref{sec:03-3}で説明しました）。同じ研究について複数の報告（本論文、続報、学会抄録）があるとき、報告は別々に数え、研究は 1 つと数えます。架空の例を図 11.1 に示します。

\begin{kkcenter}
\begin{tikzpicture}[
  fbox/.style={kbox, align=flush center, text width=46mm, font=\scriptsize, minimum height=9mm},
  xbox/.style={kplain, text width=50mm, font=\scriptsize, align=left, minimum height=9mm},
  node distance=5mm]
  \node[fbox] (a) {データベースの検索で得た記録\\2,315 件};
  \node[fbox, below=of a] (b) {重複を除いた記録\\1,708 件};
  \node[fbox, below=of b] (c) {題と抄録を選別した記録\\1,708 件};
  \node[fbox, below=of c] (d) {本文を確認した報告\\88 件};
  \node[fbox, below=12mm of d, fill=kkaccent!25] (e) {組み入れた研究\\10 件（3,626 人）\\報告 13 件};
  \node[xbox, right=6mm of b] (xb) {重複として除いた記録 607 件};
  \node[xbox, right=6mm of c] (xc) {題と抄録で除外 1,620 件};
  \node[xbox, right=6mm of d, yshift=-6mm, text width=50mm] (xd) {本文で除外 75 件\\
    \hspace*{1em}2×2 表が作れない 24\\
    \hspace*{1em}参照基準が血液培養でない 17\\
    \hspace*{1em}症例対照デザイン 12\\
    \hspace*{1em}対象が小児のみ 9\\
    \hspace*{1em}対象が発熱患者でない 8\\
    \hspace*{1em}学会抄録のみで情報が不足 5};
  \draw[kar] (a) -- (b); \draw[kar] (b) -- (c); \draw[kar] (c) -- (d); \draw[kar] (d) -- (e);
  \draw[kar] (b) -- (xb); \draw[kar] (c) -- (xc); \draw[kar] (d.east) -- (d.east -| xd.west);
\end{tikzpicture}
\kkcap{図 11.1　架空の PCT のレビューのフロー図（数は例）}
\end{kkcenter}

除外の理由は、本文を確認した段階では件数とともに報告します。抄録の段階の除外は件数だけでかまいません。組み入れた研究には、研究の数と人数に加えて、その研究にまとめた報告の数も書きます。例題の 10 研究の合計は 3,626 人で、そのうち血液培養が陽性だったのは 556 人です。

同じ研究の重複報告は、除外の理由には数えず、報告を研究にまとめる作業として扱います。同じ患者集団が、学会抄録、本論文、別の解析を加えた続報として何度も発表されることがあります。同じ患者を 2 回数えると、その研究の重みが不当に大きくなります。著者名、施設、登録期間、人数を照らし合わせて、同じ集団かどうかを判断し、1 つの研究にまとめます。解析に使うデータは、情報が最も多い報告から取ります。例題では、本文を確認した 88 件の報告のうち 75 件を除外し、残った 13 件の報告が 10 研究に対応しました。3 件の報告が、ほかの報告と同じ研究のものだったためです。

\section{この章のまとめ}
\label{sec:11-7}

検査のレビューの問いは、対象、対象検査、参照基準、標的疾患と、検査を使う場面と役割で立てます。検索ではデザインのフィルターをかけず、2 人が独立に抄録を選別し、迷った記録は本文まで残します。症例対照デザイン、参照基準が問いと違う研究、検査の結果を見て正解を決めた研究は、精度を偏らせるので除外するか、バイアスのリスクが高いと判定します。次の第 12 章では、残した論文の本文から 2×2 表を取り出し、QUADAS-2 でバイアスのリスクを評価します。
